\documentclass[11pt, a4paper, oneside]{article}
\usepackage[ngerman]{babel}
\usepackage{worksheet}

\begin{document}
	\author{L. Bung}
	\title{Webentwicklung: \hspace{10cm} HTML-Formulare}
	\subject{SAE}
	\class{E3FI}
	\maketitle
	
	Neben der Präsentation von Informationen (wie zum Beispiel Textabschnitten oder Bildern) ist es häufig notwendig, Informationen vom Nutzer abzufragen.
	Beispiele dafür sind die Anmeldung zu einem Newsletter, die Abfrage einer Lieferadresse oder die Registrierung eines Benutzerkontos.
	
	In HTML werden dafür Formulare eingesetzt, welche aus speziellen Tags bestehen.
	
	
	\partnertask{Tags für HTML-Formulare}
	
	\begin{enumerate}[label=\alph*)]
		\item Recherchieren Sie, welche HTML-Tags für den Aufbau von Formularen verwendet werden.
		\item Finden Sie heraus, was das HTML-Attribut \texttt{type} bei \texttt{<input>}-Tags bewirkt.
		\item Wie sieht es mit dem \texttt{name}-Attribut bei \texttt{<input>}-Tags aus?
		\item Was ist der korrekte Weg, um zu einem Eingabefeld einen Beschreibungstext hinzuzufügen?
	\end{enumerate}
	
	\lines[14cm]
	
	
	\singletask{Implementierung}
	
	Erstellen Sie HTML-Formulare für die in der letzten Stunde erstellten Webseite Ihres Unternehmens.
	
	Es sollen folgende Seiten mit Formularen entwickelt werden:
	
	\begin{enumerate}[label=\alph*)]
		\item Eine Anmeldeseite und -formular zu einem E-Mail-Newsletter
		\item Eine Seite zur Registrierung eines neuen Kundenkontos
		\item Ein Bestellformular, bei der Liefer- und Rechnungsadresse abgefragt werden sollen
	\end{enumerate}
	
	
	\task{Attribute von Formularelementen}
	
	\begin{enumerate}[label=\alph*)]
		\item Recherchieren Sie, welche weiteren Attribute es für die verschiedenen Formularelemente gibt und was diese bewirken.
		\item Wohin werden die eingegebenen Daten beim Absenden des Formulares gesendet?
	\end{enumerate}
	
	\lines[8cm]
	
	
	\bonustask{Weitere Formularelemente}
	
	\begin{enumerate}[label=\alph*)]
		\item Recherchieren Sie weitere Formularelemente, wie beispielsweise Farbwähler oder Zeitangaben.
		\item Erstellen Sie ein (möglichst sinnvolles...) Formular, in dem diese Elemente verwendet werden.
		\item Sehen Sie sich ein Formular an, das ein Mitschüler erstellt hat.
		Versuchen Sie anschließend, es nachzubauen, ohne den zugehörigen Code zu sehen.
	\end{enumerate}
	
\end{document}